% !TEX root = IE3_expD_prelab.tex
\documentclass[lualatex,a4paper,10pt]{jlreq}
% Shared LuaLaTeX preamble.
\usepackage{luatexja-fontspec}
\setmainfont{TeX Gyre Termes}
\setsansfont{TeX Gyre Heros}
\setmonofont{Latin Modern Mono}
\setmainjfont[BoldFont=HaranoAjiGothic-Medium]{HaranoAjiMincho-Regular}
\setsansjfont[BoldFont=HaranoAjiGothic-Bold]{HaranoAjiGothic-Medium}
\usepackage[top=22mm,bottom=22mm,left=23mm,right=23mm]{geometry}
\usepackage{amsmath,amssymb,booktabs,array,longtable,tabularx}
\usepackage{graphicx,xcolor,tikz,circuitikz}
\usetikzlibrary{arrows.meta,positioning,calc}
\usepackage{enumitem,fancyhdr,listings,needspace}
\usepackage[unicode,colorlinks=true,linkcolor=labblue,urlcolor=labteal]{hyperref}
\usepackage{bookmark}
\definecolor{labblue}{RGB}{30,70,112}
\definecolor{labteal}{RGB}{0,100,105}
\definecolor{labgray}{gray}{0.95}
\ModifyHeading{section}{font={\large\sffamily\bfseries\color{labblue}}}
\ModifyHeading{subsection}{font={\normalsize\sffamily\bfseries\color{labteal}}}
\setlength{\parindent}{1\zw}
\setlength{\parskip}{3pt}
\setlength{\emergencystretch}{2em}
\renewcommand{\arraystretch}{1.35}
\setlist{itemsep=3pt,topsep=4pt,leftmargin=2\zw}
\newcolumntype{Y}{>{\raggedright\arraybackslash}X}
\newcolumntype{C}[1]{>{\centering\arraybackslash}p{#1}}
\pagestyle{fancy}
\fancyhf{}
\fancyhead[L]{\small\color{labblue} IE3 後期・電子工学実験}
\fancyhead[R]{\small テーマ\LabCode}
\fancyfoot[C]{\thepage}
\setlength{\headheight}{16pt}
\DeclareOldFontCommand{\rm}{\normalfont\rmfamily}{\mathrm}
\lstset{basicstyle=\small\ttfamily,columns=fullflexible,keepspaces=true,
 breaklines=true,frame=single,backgroundcolor=\color{labgray},
 numbers=left,numberstyle=\tiny,showstringspaces=false,aboveskip=8pt,belowskip=8pt}
\newcommand{\blank}{\rule{22mm}{0.3pt}}
\newcommand{\answerline}{\par\noindent\rule{\linewidth}{0.25pt}\par}
\newcommand{\writing}[1]{\par\Needspace{\dimexpr#1+5mm\relax}\noindent%
 \fcolorbox{labblue!35}{white}{\parbox[c][#1][t]{\dimexpr\linewidth-2\fboxsep-2\fboxrule\relax}{\vspace{2mm}\noindent\strut\hfill\par}}\par}
\newcommand{\hint}[1]{\par\smallskip\noindent{\sffamily\bfseries\color{labteal}記述の視点：}#1\par}
\newcommand{\note}[1]{\par\smallskip\noindent\fcolorbox{labblue!45}{labblue!5}{%
 \parbox{\dimexpr\linewidth-2\fboxsep-2\fboxrule\relax}{#1}}\par\smallskip}
\newcommand{\eqerror}{\varepsilon=\frac{x_{\mathrm{meas}}-x_{\mathrm{th}}}{x_{\mathrm{th}}}\times100\ \%}
\newcommand{\LabSource}{徳山工業高等専門学校 情報電子工学科，
 『2026年度 電子工学実験（3年後期）テキスト』，テーマ\LabCode，
 \expandafter\path\expandafter{\SourcePdf}。}
\newcommand{\reporttitle}{
 \noindent{\small\sffamily\color{labteal}実験報告書・記入ガイド}\par
 \noindent{\LARGE\sffamily\bfseries \LabCode\quad\LabTitle}\par\smallskip
 \noindent 氏名：\blank\quad 班：\blank\quad 実験日：\blank\par
 \note{目的・原理・手順を確認し、実測値と自分の考察を記入する。
 考察のガイドは、検討する事象・使う測定結果・記述の方針を示す。}
}
\newcommand{\prelabtitle}{
 \noindent{\small\sffamily\color{labteal}事前学習・前提知識プリント}\par
 \noindent{\LARGE\sffamily\bfseries \LabCode\quad\LabTitle}\par\smallskip
 \noindent 氏名：\blank\quad 班：\blank\quad 予習日：\blank\par
 \note{用語、回路の動き、計算、測定表への記入の順に読む。
 例題の値は説明用であり、実験結果ではない。}
}
\newcommand{\tocrow}[3]{\hyperref[#1]{#2}& #3 &\pageref*{#1}\\[2mm]}
\newcommand{\documentmap}[1]{
 \section*{目次と概要}
 \noindent 青色の項目をクリックすると該当箇所へ移動する。\par
 {\small\begin{longtable}{@{}p{0.29\linewidth}p{0.61\linewidth}r@{}}
 \toprule {\color{labblue}\bfseries 項目} & {\color{labblue}\bfseries 読むと分かること} & 頁\\\midrule
 \endhead #1\bottomrule
 \end{longtable}}
 \clearpage
}
\newenvironment{terms}{\par\smallskip\begin{longtable}{@{}p{0.23\linewidth}p{0.73\linewidth}@{}}
 \toprule {\color{labteal}\bfseries 用語}&{\color{labteal}\bfseries 意味・回路での役割}\\\midrule\endhead}
 {\bottomrule\end{longtable}}
\newcommand{\term}[2]{\textbf{#1}&#2\\[2mm]}
\newcommand{\guidepoint}[4]{
 \Needspace{32mm}\subsection{#1}
 \noindent{\bfseries\color{labteal}考える事象：}#2\par
 \noindent{\bfseries\color{labteal}参照する結果：}#3\par
 \noindent{\bfseries\color{labteal}記述の方針：}#4\par
}
\newcommand{\reportclosing}[1]{
 \clearpage\section{結論のまとめ方}\label{sec:closing}
 \hint{#1}
 \ConclusionGuide
 \begin{enumerate}
 \item 目的に対応する最も重要な実測結果を、測定条件と数値を添えて述べる。
 \item 原理と一致した点・一致しなかった点を示す。原因は確認済みの事実と推測を分ける。
 \item 実施範囲で言えることと、未確認の条件を一文ずつまとめる。
 \end{enumerate}
 \note{書き出しの型：「○○の条件で△△を測定したところ、□□となった。
 これは◇◇と整合する。ただし××は確認しておらず、一般化には追加測定が必要である。」
 空欄は自分の実測値と根拠で埋める。}
 \noindent 結論（実験後に記入）：\writing{30mm}
 \noindent 改善案・次に確認したい条件：\writing{16mm}
 \section*{参考文献}\label{sec:references}
 \begin{enumerate}
 \item \LabSource
 \item 使用したデータシート・書籍・ウェブ資料（著者、題名、該当箇所、URL、参照日）を追加する。
 \end{enumerate}
}
\newcommand{\prelabclosing}{
 \section*{準備・出典}\label{sec:prelabsource}
 \begin{itemize}[nosep]
 \item 資料、筆記具、記録用紙、電卓と、実験条件・機器設定の記録欄。
 \item 確認問題の解答と、分からない用語・回路点への具体的な質問。
 \end{itemize}
 {\small\noindent\LabSource\par}
}
\newcommand{\simpleflow}[1]{
 \begin{center}
 \begin{tikzpicture}[node distance=5mm and 5mm,
 every node/.style={font=\small},box/.style={draw=labblue,fill=labblue!4,rounded corners=1mm,
 align=center,minimum height=10mm,text width=30mm}]
 #1
 \end{tikzpicture}
 \end{center}
}

\newcommand{\LabCode}{D}
\newcommand{\LabTitle}{LC発振回路}
\newcommand{\SourcePdf}{IE3_expD_A4.pdf}
\begin{document}
\prelabtitle
\documentmap{
\tocrow{sec:auto1}{1 用語を一つずつ理解する}{言葉の意味、単位、回路や測定での役割を確認する。}
\tocrow{sec:auto2}{2 Clapp回路で信号が一周する様子}{部品と信号の経路を追い、状態が変化する理由を理解する。}
\tocrow{sec:auto3}{3 発振と帰還}{起動と定常発振の利得・位相条件を理解する。}
\tocrow{sec:auto4}{4 LC共振の計算}{三容量の合成から理論周波数を求める。}
\tocrow{sec:auto5}{5 測定の前提}{この項目の仕組みと実験で確認すべき点を整理する。}
\tocrow{sec:auto6}{6 Fパラメータの整理}{電流方向を定義して接続条件を立てる。}
\tocrow{sec:auto7}{7 測定からレポートへ：周波数の比較}{実測値から計算・作図・記述へ進む順序を例題で練習する。}
\tocrow{sec:auto8}{8 確認問題}{自分で計算・説明して、実験に必要な理解を確かめる。}
\tocrow{sec:auto9}{解答の目安}{数値と理由を照合し、単位や論理の取り違えを見直す。}
\tocrow{sec:prelabsource}{準備・出典}{持参物と資料を確認し、具体的な質問を準備する。}
}

\section{用語を一つずつ理解する}\label{sec:auto1}
\begin{terms}
\term{発振}{外から周期的な信号を入れずに、電源のエネルギーで周期信号を生み出すこと。出力が大きくなるだけの増幅と区別する。}
\term{帰還・ループ利得}{出力の一部を入力へ戻す経路が帰還。増幅と帰還を一周した倍率がループ利得$A\beta$。大きさだけでなく位相も重要である。}
\term{正帰還}{戻った信号が元の変化を強める関係。反転増幅器が含まれていても、ループ全体の位相で判断する。}
\term{起動と定常状態}{起動時は微小な雑音等から振幅が増える。定常時は素子の非線形性等が振幅を抑えて一定になる。$|A\beta|=1$だけで必ず起動するとは言えない。}
\term{共振・合成容量}{LとCがエネルギーを交換する。Clapp形の理想モデルでは三容量の直列合成を用いる。容量の単純な和ではない。}
\term{寄生容量・負荷}{配線、トランジスタ、プローブにも小容量がある。意図したCが小さい条件では、その影響が相対的に大きくなりやすい。}
\term{リサージュ図}{二信号を水平・垂直軸へ入れたXY図。周波数比と位相で形が決まる。時間波形の横軸とは異なり、単位は両軸とも電圧である。}
\end{terms}
\section{Clapp回路で信号が一周する様子}\label{sec:auto2}
資料図D-4のトランジスタは、直流電源から動作エネルギーを受け取る。$R_1,R_2$等は直流の動作点を設定し、LC部分は交流の特定周波数を選びやすくする。
ベース側の小さな変化が増幅され、容量分割を含む経路で戻ってくる。戻る信号が位相と大きさの条件を満たすと、同じ周波数の変化が維持される。
\subsection{調整抵抗と観測点}
資料で調整するのは$R_F$であり、$R_2$は固定する。$R_E$と名前を取り違えない。
最初はL=1 mH、$C_1=C_2=C_S=0.01\,\mu\mathrm{F}$で動作を確認する。
出力端子の結合コンデンサは直流成分の扱いに関係し、共振容量三つと役割が異なる。
本プリントの$C_3$は資料の$C_S$を指す。
\subsection{三容量の計算を先に行う理由}
$C_1=C_2=C_3=C$なら、$1/C_{\rm eq}=3/C$なので$C_{\rm eq}=C/3$。
Clapp形で$C_3$だけ十分小さくした設計では$C_{\rm eq}$を主に$C_3$で決められる。
実験は等容量の組合せであり、その設計条件と同一ではない。

\section{発振と帰還}\label{sec:auto3}
発振器は直流電源からエネルギーを受け取り、外部からの周期的入力なしに周期信号を生じる。
ループ伝達関数を$A(j\omega)\beta(j\omega)$とすると、持続発振の必要条件は
\[
 |A\beta|=1,\qquad \arg(A\beta)=2\pi n.
\]
これは起動や振幅安定化まで保証する十分条件ではない。起動時の利得、飽和、バイアスも必要になる。
\simpleflow{
\node[box] (a) {増幅回路$A$};
\node[box,right=15mm of a] (b) {出力};
\node[box,below=of b] (c) {帰還回路$\beta$};
\draw[-{Stealth}] (a)--(b);\draw[-{Stealth}] (b)--(c);
\draw[-{Stealth}] (c)-| (a);
}
\section{LC共振の計算}\label{sec:auto4}
\[
 C_{\rm eq}=\left(C_1^{-1}+C_2^{-1}+C_3^{-1}\right)^{-1},
 \qquad f_0=(2\pi\sqrt{LC_{\rm eq}})^{-1}.
\]
\textbf{例題：}$L=1\,\mathrm{mH}$、各$C=0.01\,\mu\mathrm{F}$なら
$C_{\rm eq}=3.333\,\mathrm{nF}$、$f_0\simeq87.2\,\mathrm{kHz}$となる。
Lを4倍にすると周波数は半分となる。
$C_3$のみを十分小さくすると$C_{\rm eq}\simeq C_3$となり、他の容量や寄生容量の影響を抑えやすい。
\clearpage
\section{測定の前提}\label{sec:auto5}
リサージュ図は2信号の周波数比と位相関係を表す。まず同じ周波数に合わせると閉じた楕円等が静止する。複雑な図では周波数の上下関係を取り違えないよう、既知の基準と比較する。
周波数計も、トリガレベルや波形のひずみで誤計数し得る。
\section{Fパラメータの整理}\label{sec:auto6}
資料の電流方向の定義に従って
\[
 \begin{bmatrix}V_i\\I_i\end{bmatrix}
 =F\begin{bmatrix}V_o\\I_o\end{bmatrix}
\]
と書く。増幅回路と帰還回路の接続条件を同じ方向規約で整理すると、非零の電圧・電流が存在する条件は係数行列の行列式が零になることである。
自分の回路図に電流の矢印を描いてから式を立てる。
\section{測定からレポートへ：周波数の比較}\label{sec:auto7}
\subsection{単位を直して九組を計算する}
例として$L=470\,\mu\mathrm{H}=4.70\times10^{-4}\,\mathrm{H}$、各$C=0.01\,\mu\mathrm{F}=10^{-8}\,\mathrm{F}$。
合成容量は$3.333\times10^{-9}\,\mathrm{F}$なので、理論周波数は約127 kHzとなる。
mH、$\mu$H、$\mu$Fを数字だけで代入すると桁が大きく変わる。表には各Cと合成Cを区別して示す。
\subsection{二種類の測定をそろえる}
XY図が1対1で静止したら、基準発振器の周波数を周波数計で読んで記録する。
比が異なる図では、既知の基準側と実験側を区別して周波数比を確認する。
次に実験回路へ周波数計を直接接続し、出力波形が変化していないか確認する。
\subsection{誤差とグラフ}
仮に理論100 kHz、実測97 kHzなら相対誤差は$-3\,\%$。符号は低い方向への差を示す。
両対数グラフでは横軸を理論周波数、縦軸を実測周波数とし、一致線$y=x$を引く。
L、Cの条件を点の記号等で区別すると、高周波側で差が増える傾向を調べやすい。
「測定誤差」とひとまとめにせず、寄生容量・部品の実値・計器負荷・読取りを候補ごとに検討する。

\clearpage
\clearpage
\section{確認問題}\label{sec:auto8}
\begin{enumerate}
\item 同じ値Cのコンデンサ3個の直列合成容量を求める。\answerline
\item 上記例題で各容量を10倍にすると周波数は何倍になるか。\answerline
\item 理論$100\,\mathrm{kHz}$、実測$95\,\mathrm{kHz}$の相対誤差を求める。\answerline
\item 出力端子にプローブを接続すると周波数が変化する理由を挙げる。\answerline
\end{enumerate}
\section*{解答の目安}\label{sec:auto9}
(1) $C/3$。(2) $1/\sqrt{10}\simeq0.316$倍。(3) $-5\,\%$。
(4) プローブの入力容量・抵抗が共振条件や負荷を変えるため。
\prelabclosing
\end{document}
